% GENERATED FILE. Edit canonical lessons, then run npm run generate:books.
% canonical-source-hash: fnv1a64:a7525c254133ef61
% canonical-lessons: MR-C398-tikitghar, MR-C398-pas, MR-C398-partava, MR-C398-gantavya, MR-C398-rakhun-thevne

\chapter{At the Ticket Office}
\label{ch:mr-at-the-ticket-office}
\hlchaptermodality{398}

\begin{chapteropening}
\textbf{By the end of this chapter:} \emph{I can talk about a ticket office, a pass, a refund and where I am going.}

Complete the last lesson of chapter 398: I can talk about a ticket office, a pass, a refund and where I am going.
\end{chapteropening}

\section[\texorpdfstring{\mr{तिकीटघर} (tikīṭghar)}{तिकीटघर (tikīṭghar)}]{\mr{तिकीटघर} (tikīṭghar) — a ticket office}

\label{lesson:MR-C398-tikitghar}



\begin{quote}
Before the new one: say the Marathi for to be stalled, then the Marathi for to set off.
\end{quote}

\subsection*{You'll want to know: \mr{तिकीटघर}}
\textbf{\mr{तिकीटघर}} — \emph{tikīṭghar} — \textquotedblleft{}a ticket office\textquotedblright{}.

\begin{quote}
\mr{तिकीटघरासमोर} \mr{मोठी} \mr{रांग} \mr{होती}.
\end{quote}

\begin{quote}
\emph{tikīṭgharāsamor moṭhī rāṅg hotī.}
\end{quote}

\begin{quote}
There was a long queue in front of the ticket office.
\end{quote}

Neuter.

\mr{घर} at the end of a compound is a room or building for one job: \mr{तिकीटघर}, a ticket house.

\begin{cousinweb}
\textbf{\mr{तिकीटघर}} (\emph{tikīṭghar}) is a Marathi compound of \textbf{\mr{तिकीट}}, a loan from English \emph{ticket}, and \textbf{\mr{घर}}, \textquotedblleft{}house\textquotedblright{}, inherited from Sanskrit \emph{ghara} (CDIAL 4428).
\end{cousinweb}

\begin{grammarlens}[title={What you've built}]
One more word for getting around.
\end{grammarlens}

\subsection*{Your turn}
\begin{itemize}
\raggedright
  \item \emph{Say it:} \emph{tikīṭghar}
  \item \emph{Say it:} \emph{tikīṭghar}, once more
  \item \emph{Say it:} \emph{tikīṭgharāsamor moṭhī rāṅg hotī.} — There was a long queue in front of the ticket office.
  \item \emph{Recall:} say the Marathi for to be stalled, then the Marathi for to set off, then say \emph{tikīṭghar} again
\end{itemize}

\begin{tcolorbox}[breakable,colback=teal!4,colframe=teal!35!black,title={Before you move on}]
What does \mr{तिकीटघर} mean? (\textquotedblleft{}A ticket office\textquotedblright{}.) Say it once more.
\end{tcolorbox}
\section[\texorpdfstring{\mr{पास} (pās)}{पास (pās)}]{\mr{पास} (pās) — a travel pass; a season ticket}

\label{lesson:MR-C398-pas}



\begin{quote}
Before the new one: say the Marathi for to set off, then the Marathi for a ticket office.
\end{quote}

\subsection*{You'll want to know: \mr{पास}}
\textbf{\mr{पास}} — \emph{pās} — \textquotedblleft{}a travel pass; a season ticket\textquotedblright{}.

\begin{quote}
\mr{अरुणकडे} \mr{बसचा} \mr{महिन्याचा} \mr{पास} \mr{आहे}.
\end{quote}

\begin{quote}
\emph{aruṇkaḍe bascā mahinyācā pās āhe.}
\end{quote}

\begin{quote}
Arun has a monthly bus pass.
\end{quote}

Masculine.

A pass lasts a week, a month or a year: \mr{आठवड्याचा} \mr{पास}, a weekly pass.

\begin{cousinweb}
\textbf{\mr{पास}} (\emph{pās}) is a loan from English \emph{pass}, the card that lets you travel again and again for one price.
\end{cousinweb}

\begin{grammarlens}[title={What you've built}]
One more word for getting around.
\end{grammarlens}

\subsection*{Your turn}
\begin{itemize}
\raggedright
  \item \emph{Say it:} \emph{pās}
  \item \emph{Say it:} \emph{pās}, once more
  \item \emph{Say it:} \emph{aruṇkaḍe bascā mahinyācā pās āhe.} — Arun has a monthly bus pass.
  \item \emph{Recall:} say the Marathi for to set off, then the Marathi for a ticket office, then say \emph{pās} again
\end{itemize}

\begin{tcolorbox}[breakable,colback=teal!4,colframe=teal!35!black,title={Before you move on}]
What does \mr{पास} mean? (\textquotedblleft{}A travel pass; a season ticket\textquotedblright{}.) Say it once more.
\end{tcolorbox}
\section[\texorpdfstring{\mr{परतावा} (partāvā)}{परतावा (partāvā)}]{\mr{परतावा} (partāvā) — a refund}

\label{lesson:MR-C398-partava}



\begin{quote}
Before the new one: say the Marathi for a ticket office, then the Marathi for a travel pass.
\end{quote}

\subsection*{You'll want to know: \mr{परतावा}}
\textbf{\mr{परतावा}} — \emph{partāvā} — \textquotedblleft{}a refund\textquotedblright{}.

\begin{quote}
\mr{मला} \mr{तिकिटाचा} \mr{पूर्ण} \mr{परतावा} \mr{मिळाला}.
\end{quote}

\begin{quote}
\emph{malā tikiṭācā pūrṇa partāvā miḷālā.}
\end{quote}

\begin{quote}
I got a full refund for the ticket.
\end{quote}

Masculine.

Built on \mr{परत}, back: a refund is money that comes back to you.

\begin{cousinweb}
\textbf{\mr{परतावा}} (\emph{partāvā}) is a Marathi derivation built on \textbf{\mr{परत}}, \textquotedblleft{}back, again\textquotedblright{}, from the family of Sanskrit \emph{parivartate}, \textquotedblleft{}turns round\textquotedblright{}, where CDIAL lists Marathi \emph{paratṇẽ}, \textquotedblleft{}to return\textquotedblright{} (CDIAL 7872).
\end{cousinweb}

\begin{grammarlens}[title={What you've built}]
One more word for getting around.
\end{grammarlens}

\subsection*{Your turn}
\begin{itemize}
\raggedright
  \item \emph{Say it:} \emph{partāvā}
  \item \emph{Say it:} \emph{partāvā}, once more
  \item \emph{Say it:} \emph{malā tikiṭācā pūrṇa partāvā miḷālā.} — I got a full refund for the ticket.
  \item \emph{Recall:} say the Marathi for a ticket office, then the Marathi for a travel pass, then say \emph{partāvā} again
\end{itemize}

\begin{tcolorbox}[breakable,colback=teal!4,colframe=teal!35!black,title={Before you move on}]
What does \mr{परतावा} mean? (\textquotedblleft{}A refund\textquotedblright{}.) Say it once more.
\end{tcolorbox}
\section[\texorpdfstring{\mr{गंतव्य} (gantavya)}{गंतव्य (gantavya)}]{\mr{गंतव्य} (gantavya) — a destination}

\label{lesson:MR-C398-gantavya}



\begin{quote}
Before the new one: say the Marathi for a travel pass, then the Marathi for a refund.
\end{quote}

\subsection*{You'll want to know: \mr{गंतव्य}}
\textbf{\mr{गंतव्य}} — \emph{gantavya} — \textquotedblleft{}a destination\textquotedblright{}.

\begin{quote}
\mr{आमचं} \mr{गंतव्य} \mr{अजून} \mr{खूप} \mr{दूर} \mr{आहे}.
\end{quote}

\begin{quote}
\emph{āmcaṁ gantavya ajūn khūp dūr āhe.}
\end{quote}

\begin{quote}
Our destination is still very far away.
\end{quote}

Neuter.

A formal word; in speech people say \mr{जिथे} \mr{जायचं} \mr{आहे} \mr{ते} \mr{ठिकाण}.

\begin{cousinweb}
\textbf{\mr{गंतव्य}} (\emph{gantavya}) is a tatsama, taken from Sanskrit \emph{gantavya}, \textquotedblleft{}to be gone to\textquotedblright{}, a form of the verb \emph{gam}, \textquotedblleft{}to go\textquotedblright{} (Monier-Williams).
\end{cousinweb}

\begin{grammarlens}[title={What you've built}]
One more word for getting around.
\end{grammarlens}

\subsection*{Your turn}
\begin{itemize}
\raggedright
  \item \emph{Say it:} \emph{gantavya}
  \item \emph{Say it:} \emph{gantavya}, once more
  \item \emph{Say it:} \emph{āmcaṁ gantavya ajūn khūp dūr āhe.} — Our destination is still very far away.
  \item \emph{Recall:} say the Marathi for a travel pass, then the Marathi for a refund, then say \emph{gantavya} again
\end{itemize}

\begin{tcolorbox}[breakable,colback=teal!4,colframe=teal!35!black,title={Before you move on}]
What does \mr{गंतव्य} mean? (\textquotedblleft{}A destination\textquotedblright{}.) Say it once more.
\end{tcolorbox}
\section[\texorpdfstring{\mr{राखून} \mr{ठेवणे} (rākhūn ṭhevṇe)}{राखून ठेवणे (rākhūn ṭhevṇe)}]{\mr{राखून} \mr{ठेवणे} (rākhūn ṭhevṇe) — to set aside; to keep back}

\label{lesson:MR-C398-rakhun-thevne}



\begin{quote}
Before the new one: say the Marathi for a refund, then the Marathi for a destination.
\end{quote}

\subsection*{You'll want to know: \mr{राखून} \mr{ठेवणे}}
\textbf{\mr{राखून} \mr{ठेवणे}} — \emph{rākhūn ṭhevṇe} — \textquotedblleft{}to set aside; to keep back\textquotedblright{}.

\begin{quote}
\mr{मी} \mr{तुझ्यासाठी} \mr{एक} \mr{जागा} \mr{राखून} \mr{ठेवली}.
\end{quote}

\begin{quote}
\emph{mī tujhyāsāṭhī ek jāgā rākhūn ṭhevlī.}
\end{quote}

\begin{quote}
I kept a seat aside for you.
\end{quote}

Past, with -\mr{ने} (he …): \mr{त्याने} \mr{राखून} \mr{ठेवलं}, \emph{tyāne rākhūn ṭhevlaṁ}.

\mr{राखणे} is to guard or keep; with \mr{ठेवणे} it means to hold something back for later. The verb agrees with the thing kept: \mr{जागा} \mr{राखून} \mr{ठेवली}.

\begin{cousinweb}
\textbf{\mr{राखून} \mr{ठेवणे}} (\emph{rākhūn ṭhevṇe}) joins \textbf{\mr{राखून}}, the -\mr{ऊन} form of \textbf{\mr{राखणे}}, \textquotedblleft{}to keep, to guard\textquotedblright{}, inherited from Sanskrit \emph{rakṣati}, \textquotedblleft{}guards\textquotedblright{} (CDIAL 10547), to \textbf{\mr{ठेवणे}}, \textquotedblleft{}to put, to keep\textquotedblright{}.
\end{cousinweb}

\begin{grammarlens}[title={What you've built}]
One more word for getting around.
\end{grammarlens}

\subsection*{Your turn}
\begin{itemize}
\raggedright
  \item \emph{Say it:} \emph{rākhūn ṭhevṇe}
  \item \emph{Say it:} \emph{rākhūn ṭhevṇe}, once more
  \item \emph{Say it:} \emph{mī tujhyāsāṭhī ek jāgā rākhūn ṭhevlī.} — I kept a seat aside for you.
  \item \emph{Recall:} say the Marathi for a refund, then the Marathi for a destination, then say \emph{rākhūn ṭhevṇe} again
\end{itemize}

\begin{tcolorbox}[breakable,colback=teal!4,colframe=teal!35!black,title={Before you move on}]
What does \mr{राखून} \mr{ठेवणे} mean? (\textquotedblleft{}To set aside; to keep back\textquotedblright{}.) Say it once more.
\end{tcolorbox}
